\subsection*{完整示例：逐元素加法}

下面的程序对两个形状为 \co@[4, 8]@ 的 32 位整数矩阵逐元素求和：

\begin{lstlisting}[language=croqtile,caption={完整的逐元素加法程序}]
__co__ s32 [4, 8] ele_add(s32 [4, 8] lhs, s32 [4, 8] rhs) {
  s32 [lhs.span] output;
  parallel {i, j} by [4, 8]
    output.at(i, j) = lhs.at(i, j) + rhs.at(i, j);
  return output;
}
int main() {
  auto lhs = croq::make_spandata<croq::s32>(4, 8);
  auto rhs = croq::make_spandata<croq::s32>(4, 8);
  lhs.fill_random(-10, 10);
  rhs.fill_random(-10, 10);
  auto res = ele_add(lhs.view(), rhs.view());
  for (int i = 0; i < 4; ++i)
    for (int j = 0; j < 8; ++j)
      croq_assert(lhs[i][j] + rhs[i][j] == res[i][j],
                          "values are not equal.");
  std::cout << "Test Passed\n" << std::endl;
}
\end{lstlisting}

将代码保存为 \co@ele_add.co@，然后编译并运行：

\begin{lstlisting}[language=croqtile,caption={编译并运行加法示例}]
croq ele_add.co -o ele_add
./ele_add
\end{lstlisting}

如果输出 \co@Test Passed@，说明程序计算出的结果通过了检查。第 2 章会进一步拆解这个示例，说明函数声明、形状和并行语句各自的作用。
